\chapter{Sleep Study}
\section*{What Is This Test?} A sleep study records breathing, oxygen, brain waves, heart rhythm, movement, and sleep stages during sleep.\section*{Alternative Names} Polysomnography; overnight sleep study; home sleep apnea test.\section*{Principle} Sensors detect airflow, respiratory effort, oxygen, EEG, muscle activity, position, and heart rate to identify sleep and breathing events.\section*{Why This Test Is Ordered} It evaluates suspected sleep apnea, unusual movements, narcolepsy, parasomnias, or persistent daytime sleepiness.\section*{Primary vs Secondary Test} Clinical sleep assessment is primary; polysomnography or home testing is selected for the suspected disorder.\section*{How This Test Is Performed} Sensors are applied before sleep in a lab or a validated home device is used overnight.\section*{What Preparation Is Needed} Follow instructions about caffeine, alcohol, naps, medicines, and hair or skin products.\section*{Understanding Results} The report may show apnea-hypopnea index, oxygen changes, arousals, sleep stages, limb movements, and rhythm findings.\section*{Follow-up Tests} Titration study, sleep specialist review, pulmonary or cardiac assessment, or repeat testing may follow.\section*{References} \begin{enumerate}\item MedlinePlus. Sleep Study.\item American Academy of Sleep Medicine. Clinical practice guidelines.\end{enumerate}\clearpage\section*{Understanding Results and Follow-up}
The laboratory report should identify the specimen, method, units, reference interval or decision limit, and whether the result is positive, negative, abnormal, or inconclusive. Reference intervals vary with age, sex, pregnancy, specimen type, assay, and laboratory; a result outside a range is not automatically a diagnosis, and a result within a range does not always exclude disease. Discuss unexpected results with the ordering clinician, who can decide whether repeat or confirmatory testing, treatment, or referral is appropriate. Do not change medicines or delay urgent care based on an isolated result.
\section*{Regulatory Status and Authorization Date}This chapter describes a test category, not a specific commercial product. FDA, CDSCO, EU IVDR, and MHRA approval or registration dates therefore cannot be assigned without the assay manufacturer, product name, instrument, and intended use. For laboratory-developed tests, an individual agency approval date may not exist. See the Preface for the interpretation of regulatory status.
\clearpage
